% =====================================================================
%  Capítulo 5. Conclusiones  (Bloque UMA obligatorio)
%  ESTADO: andamiaje. NO redactar todavía (depende de §4.6 y §4.7).
%  No introducir resultados nuevos aquí.
% =====================================================================
\chapter{Conclusiones}
\label{cap:conclusiones}

\pendienteredaccion{} Redactar tras cerrar resultados y amenazas.

\section{Respuesta a las preguntas de investigación}
\pendienteredaccion{} RQ1: precondiciones P1--P5 y transformación
canónica. RQ2: $\Delta$ CC agregado (re-extraído) y validación parcial con
SonarQube. RQ3: tasas de éxito y \emph{baseline match} por modelo en la
campaña fase~9. (Sin cifras hasta consolidar §\ref{sec:resultados}.)

\section{Contribuciones efectivas}
\pendienteredaccion{} Prototipo determinista, paquete de replicación,
protocolo experimental congelado y taxonomía de descartes.

\section{Limitaciones}
\pendienteredaccion{} Resumen de \autoref{cap:amenazas}.

\section{Trabajo futuro}
\pendienteredaccion{} Orientado a una posible continuación investigadora
(p.\,ej. tesis doctoral): extensión del catálogo de patrones
(\lstinline|else|/\lstinline|else if| simétricos, anidamiento triple
parcial, \emph{guards} sin efectos colaterales); validación SonarQube
completa sobre todo el corpus; ampliación y diversificación del corpus
real; comparativa con un segundo proveedor de LLM y con técnicas
\emph{few-shot}; e integración como \emph{plugin} de IDE o regla de
SonarQube.
